\section{The MedicalQA Dataset and Pipelines}
\label{sec:method}

% TODO: flesh out from the codebase. Skeleton mirrors the MentorQA paper
% (Sections 3-4) so the extension reads as a controlled domain transfer.

\subsection{Expert Lecture Corpus}
% TODO: describe the source lectures (expert dementia-care educator
% content, e.g., the Teepa Snow lecture series), hours, transcription
% (download_transcripts.py / preprocess.py), and any second corpus
% ("Master"). Licensing / public availability note.
\todo{Corpus description: videos, hours, transcription, preprocessing.}

\subsection{QA-Generation Architectures}
% Mirror MentorQA's four architectures; name the agents as implemented:
% SingleQA (single-agent), LLMChunking (chunking agent + QA agent),
% RAG (question generation + ChromaDB retrieval-grounded answering,
% RAG/vector_embeddings/chroma_db.py), MultiAgentChunking (architect,
% inquisitor, scorer, justifier, synthesizer — agents/agent1-5).
\todo{Describe the four pipelines and the five multi-agent roles
(architect, inquisitor, scorer, justifier, synthesizer); base model and
serving setup (Qwen weights on the cluster).}

\subsection{Grounding and Quality Control}
\todo{Segment-grounded answers, quota/scoring if applicable, PII and
safety considerations for health content.}
